\documentclass[12pt,a4paper,oneside]{article}
\usepackage[brazil]{babel}
\usepackage[utf8]{inputenc}
\usepackage{olymp}

\setcounter{problem}{0}

\begin{document}

\begin{problem}{Dividindo a conta}{divconta}{2}
 
Depois de mais uma prova fácil os estudantes saíram para festejar numa lanchonete da cidade. A conta chegou e cada um pagou certo valor. Agora precisam saber quanto cada um tem de crédito (ou débito!) para ser considerado na próxima comemoração, considerando que a conta foi dividida igualmente entre todos.

%\Notes
%
%Observações, se tiver.

\InputFile

A entrada é composta por duas linhas. A primeira contém um número inteiro $N$, a quantidade de estudantes que dividiram a conta. A segunda linha contém $N$ valores inteiros, $P_1, P_2, \dots, P_N$, que indicam o valor pago por cada estudante.
Restrições: $1 < N \leq 50$ e $0 \leq P_i \leq 100$.

\OutputFile

A saída do seu programa deve conter uma linha informando o valor de crédito (ou débito) de cada estudante arredondado para 2 casas decimais. Esses valores devem estar separados por espaço em branco, com quebra de linha após o último.

% \Notes
% É garantido que sempre existe uma divisão que satisfaz as condições dos reinos.

\Example

\begin{example}
\exmp{
5
20 10 15 30 25
}{
0.00 -10.00 -5.00 10.00 5.00
}%
\end{example}

\begin{example}
\exmp{
4
10 10 10 0
}{
2.50 2.50 2.50 -7.50
}%
\end{example}

\begin{example}
\exmp{
3
30 30 40
}{
-3.33 -3.33 6.67
}%
\end{example}

% \small \textit{Obs.: baseada na questão ``Guerra por território'' da OBI 2012} 

%Comentários sobre o exemplo, se necessário.
\small \textit{Neo primeiro exemplo, o total da conta foi 100, então cada um dos 5 estudantes deveria pagar 20. Assim, o primeiro não tem crédito nem débito, o segundo e o terceiro têm débito e os dois últimos têm crédito.}

\end{problem}

\end{document}
